\documentclass[a4paper,10pt,twoside]{article}
\usepackage{its_conf}
\begin{document}
\setcounter{section}{0}
\setcounter{figure}{0}
\setcounter{table}{0}
\setcounter{equation}{0}
\setcounter{secnumdepth}{1}
\setcounter{secnumdepth}{1}
% ======= Начало Ваших тезисов. Все, что выше - не изменять ! ===========

\topic{Адаптивный конвейер визуализации больших графов в текстово-графическом редакторе диаграмм}

\information{Хаджинова~К.~А.}
{Белорусский государственный университет информатики и радиоэлектроники}
{Минск, Республика Беларусь}
{kseniyakhadzhynava@gmail.com}

% ========= Аннотация ==========
\annotation{Рассматривается задача сохранения интерактивности текстово-графического редактора диаграмм при работе с графами из тысяч и десятков тысяч элементов. Предложен адаптивный конвейер визуализации, объединяющий выбор внутреннего представления графа, пространственный отбор видимых элементов, кластеризацию малозначимых подграфов и поэтапный рендеринг. На графах до 50\,000 узлов экспериментально сравнены три представления графа и три способа пространственного поиска. Типизированные массивы сокращают память примерно в десять раз, а равномерная сетка ускоряет поиск элемента в сотни раз по сравнению с перебором.}

\begin{multicols}{2} %НЕ УДАЛЯТЬ!!!

\section*{Введение}
В редакторах диаграмм модель процесса, данных или архитектуры фактически является атрибутированным графом: узлы соответствуют элементам, ребра~-- связям. При тысячах и десятках тысяч элементов документ открывается долго, панорамирование и масштабирование сопровождаются задержками, а определение элемента под курсором становится заметной операцией.

Стандартным способом вывода в веб-приложениях является SVG: каждый элемент становится узлом DOM, что удобно для событий и стилей, но делает массовую отрисовку дорогой. Canvas~2D рисует множество примитивов в одном элементе, однако требует собственной реализации выделения и обновления. Для больших графов применяют кластеризацию данных [1], уровни детализации, при которых плотные области графа агрегируются [2], и отображение только текущей области просмотра: так, в GraphMaps на каждом масштабе рисуются лишь элементы внутри видимой области [3], а в ASK-GraphView используется иерархическая кластеризация [4].

Цель работы~-- описать адаптивный конвейер визуализации больших графов и оценить вклад хранения графа и пространственного поиска. Ранее авторами была предложена конвейерная модель обработки диаграмм с единым графовым представлением [5]; здесь исследуются его хранение и отображение.

\section{Адаптивный конвейер}
Диаграмма задается текстом на предметно-ориентированном языке, где каждая строка описывает узел или связь. Конвейер состоит из семи этапов, показанных на рис.~\ref{fig:arch}: разбор текста, построение внутреннего представления графа, пространственный индекс, планировщик видимости, поэтапный рендерер и слой взаимодействия.

\image{
[width=\columnwidth]{fig-architecture.png}
\caption{Архитектура адаптивного конвейера\label{fig:arch}}
}

Как видно из рис.~\ref{fig:arch}, действия пользователя возвращаются к планировщику видимости. Сравниваются три представления графа: объектная модель (узлы и ребра как объекты со ссылками), списки смежности и компактное представление на типизированных массивах.

Для поиска элемента под курсором и отбора видимых элементов сравниваются перебор, равномерная сетка и quadtree. Перебор проверяет все $n$ узлов, индекс ограничивает проверку узлами нужной области. Эффективность отсечения невидимых элементов оценивается коэффициентом по формуле~(1):
$$
C_{cull} = 1 - N_{drawn}/N_{total},\eqno(1)
$$
где $N_{drawn}$ и $N_{total}$~-- число отображаемых и всех элементов.

Для поэтапного вывода узлу присваивается приоритет по формуле~(2):
$$
p(v) = 0{,}5\,deg(v) + 0{,}3\,level(v) + 0{,}2\,sel(v),\eqno(2)
$$
где $deg(v)$ и $level(v)$~-- нормализованные степень и уровень иерархии, $sel(v)$~-- признак выделения. Группы узлов с низким приоритетом (2) заменяются кластером с числом скрытых элементов.

\section{Методика эксперимента}
Измерения выполнены на отдельном стенде, не связанном с основным редактором (Node.js 24, движок V8; Intel Core i5-1035G1, 24~ГБ ОЗУ, Windows~11). Генератор строит текст на языке диаграмм для графов из 500, 2\,000, 10\,000 и 50\,000 узлов трех типов: случайных, иерархических и с узлами-хабами; начальное значение генератора фиксировано. Измеряются память модели (куча после сборки мусора), время построения модели, обхода и индекса, время поиска элемента под курсором (1\,000 точек) и отбора узлов в окне $1600\times900$ (200 окон). Каждое измерение выполняется в отдельном процессе, 11 запусков, первый отбрасывается; приводятся медианы, для трех типов графов~-- медиана по типам. Исходные данные измерений сохранены в CSV-файлах, что позволяет воспроизвести результаты.

\section{Результаты}
Разбор текста занимает около 0,55~с на 50\,000 узлов и не зависит от представления графа, поэтому остальные показатели в табл.~\ref{tab:repr} приведены без него. Типизированные массивы занимают в десять раз меньше памяти, чем объекты, строят модель в 3--6 раз быстрее и обходят граф в 10 раз быстрее.

\myTable{Представления графа\label{tab:repr}}{{\columnwidth}{|>{\raggedright\arraybackslash}p{2.5cm}|X|X|X|} \hline
Показатель & Объекты & Списки & Массивы \\ \hline
Память, МБ (50\,000) & 24,3 & 23,2 & 2,30 \\ \hline
Построение, мс (50\,000) & 200 & 98 & 31 \\ \hline
Обход, мс (50\,000) & 138 & 27 & 14 \\ \hline
}

Списки смежности экономят лишь 5\,\% памяти. Результаты пространственного поиска для 50\,000 узлов приведены в табл.~\ref{tab:idx}; индекс строится один раз при загрузке, поэтому указано и время его построения.

\myTable{Пространственный поиск, 50\,000 узлов\label{tab:idx}}{{\columnwidth}{|>{\raggedright\arraybackslash}p{2.9cm}|X|X|X|X|} \hline
Метод & Индекс, мс & Поиск, мкс & 95\,\%, мкс & Отбор, мкс \\ \hline
Перебор, объекты & 0 & 1446 & 4952 & 1789 \\ \hline
Перебор, массивы & 0 & 573 & 2086 & 775 \\ \hline
Сетка & 69 & 1,3 & 4,4 & 50 \\ \hline
Quadtree & 644 & 5,0 & 18,8 & 55 \\ \hline
}

Равномерная сетка ускоряет поиск элемента в 440 раз по сравнению с перебором по массивам, quadtree~-- в 115 раз (рис.~\ref{fig:hit}). Сетка быстрее quadtree и на графах с узлами-хабами (1,3 и 4,1~мкс против 621~мкс при переборе), а строится в 9 раз быстрее.

\image{
[width=0.85\columnwidth]{fig-hit.png}
\caption{Время поиска элемента под курсором\label{fig:hit}}
}

На рис.~\ref{fig:hit} время поиска перебором растет примерно пропорционально числу узлов, а для индексов растет слабо. Даже при 50\,000 узлов перебор по массивам занимает 0,6~мс, то есть укладывается в бюджет кадра 16,7~мс, поэтому узким местом на таких размерах ожидается отрисовка; проверка для SVG и Canvas является следующим этапом. Разбор текста (0,55~с) занимает на порядок больше времени, чем построение массивов (31~мс), поэтому перспективен инкрементальный разбор.

\section*{Заключение}
Предложен конвейер визуализации больших графов и измерено влияние представления графа и пространственного индекса. Типизированные массивы целесообразны от нескольких тысяч узлов, а при близкой к равномерной плотности элементов предпочтительнее сетка: она быстрее и дешевле в построении, чем quadtree. Ограничение работы~-- измерения выполнены в среде Node.js без отрисовки в браузере. Дальнейшая работа~-- измерение отрисовки SVG и Canvas, поэтапный вывод и инкрементальный разбор текста.

\ListReferences{
\item Козарь,~Р.~В. Методы кластеризации данных при распознавании эндоскопических снимков в задачах компьютерной медицинской диагностики~/ Р.~В.~Козарь, Н.~С.~Конойко, А.~А.~Навроцкий~// Доклады БГУИР.~-- 2023.~-- Т.~21, №~1.~-- С.~94--97.~-- DOI: 10.35596/1729-7648-2023-21-1-94-97.
\item Zinsmaier,~M. Interactive level-of-detail rendering of large graphs~/ M.~Zinsmaier, U.~Brandes, O.~Deussen, H.~Strobelt~// IEEE Trans. Vis. Comput. Graph.~-- 2012.~-- Vol.~18, №~12.~-- P.~2486--2495.
\item GraphMaps: browsing large graphs as interactive maps~/ L.~Nachmanson [et al.]~// arXiv.~-- 2015.~-- arXiv:1506.06745.
\item Abello,~J. ASK-GraphView: a large scale graph visualization system~/ J.~Abello, F.~van Ham, N.~Krishnan~// IEEE Trans. Vis. Comput. Graph.~-- 2006.~-- Vol.~12, №~5.~-- P.~669--676.
\item Ломако,~А.~В. Автоматизация преобразования изображений диаграмм в программный код~/ А.~В.~Ломако, К.~А.~Хаджинова~// Технологии передачи и обработки информации.~-- 2026.
}

\end{multicols} % НЕ УДАЛЯТЬ!

% ===== Конец материалов. Все, что ниже, не изменять !===========
\end{document}
